% ===================================================================
%  పట్టిక సంఖ్య 15 — సంత. తినేవి.
%
%  Traduction, faite en 2026, en telougou d'aujourd'hui, sur l'ido de
%  texto/io. Regles de la colonne : voir 10-pattika-01.tex. Dix-huit
%  alineas, cent cinq renvois — le plus grand nombre de renvois du
%  livret.
%
%  ET SEPT RETOURNEMENTS SEULEMENT, sur dix-huit paires. Le tableau 13
%  en coutait dix-neuf pour quatre-vingt-treize renvois. Le marche est
%  la forme extreme de ce que le tableau 14 avait montre : il ne fait
%  presque QUE des listes — le jambon, le boudin, la saucisse, le
%  lard ; les tomates, les choux-fleurs, la salade, les choux. Sur
%  cent cinq renvois, une centaine sont des membres d'enumeration, et
%  une langue a tete finale les rend dans l'ordre recu sans rien
%  faire. Les sept qui coutent sont tous des relatives — « qui vend »,
%  « qui pousse », « qui emploie », « qui amidonne » — et une seule
%  vient d'ailleurs : « livro-veturo (21) kargita per bir-bareli
%  (22) », que la liste ne donnait pas avant-hier et qu'elle donne
%  depuis que le participe passif est entre dans le motif.
%
%  LE MARCHE EST L'ENDROIT OU LE TELOUGOU EST LE PLUS CHEZ LUI DE TOUT
%  LE LIVRET. Le riz, les lentilles, l'ail, l'oignon, la banane,
%  l'ananas, le crabe, la crevette, la caille, la perdrix ne sont pas
%  des traductions : ce sont les mots de tous les jours en Andhra —
%  బియ్యం, మసూరు పప్పు, వెల్లుల్లి, ఉల్లిపాయ, అరటి, అనాస, పీత,
%  రొయ్య, పూరేడు, కౌజు. Le pays du livret et le pays de la langue
%  mangent les memes choses, et la colonne s'ecrit toute seule.
%
%  LA COUPURE, ELLE, EST TOUTE EUROPEENNE, et elle se lit d'un coup
%  d'œil : ఆర్టిచోకులు l'artichaut, కామంబేర్, గ్రుయేర్, హాలండు
%  జున్నులు les fromages, హేజల్‌నట్లు les noisettes, ఆలివ్‌లు les
%  olives, ఫీజెంట్ le faisan, క్రొసాం le croissant. Aucun de ces
%  objets n'existe dans une cuisine andhra ; ils gardent donc leur
%  nom, comme le marron chaud du tableau 5.
%
%  ET UN OISEAU DE PLUS QUE JE NE SAIS PAS NOMMER : la grive (40).
%  C'est le quatrieme du livret apres l'hirondelle, la bergeronnette
%  et l'alouette du tableau 8. Le telougou nomme la plupart des
%  oiseaux qu'il voit — కౌజు la perdrix, పూరేడు la caille, కాకి la
%  corneille — mais pas ceux qui ne passent pas au-dessus de l'Andhra.
%  On ecrit l'emprunt et on le dit ici.
%
%  DEUX METIERS DE VIANDE, ET LE MOT QU'ON N'EMPLOIE PAS. « కసాయి »
%  traduit exactement « boucher », et c'est aussi, en telougou
%  d'aujourd'hui, une insulte courante — « brute », « egorgeur ». La
%  colonne ecrit donc మాంసం అమ్మేవాడు et మాంసపు దుకాణం, la fonction
%  et le lieu. Ce n'est pas la regle des castes du tableau 6 : c'est
%  la meme raison plus simple encore — le mot juste de 1926 est
%  devenu une injure en 2026, et une traduction de 2026 ne l'ecrit
%  pas sans le vouloir.
% ===================================================================

%%K t15-apar-1 apar
\VUsaut{4.0mm}
\VUtitre{15.0pt}{60}{పట్టిక సంఖ్య 15}
\VUsaut{3.0mm}

%%K t15-tit-1 sub
\VUpk{11.4pt}{40}{సంత. --- తినేవి.}
\VUsaut{3.0mm}

%%K t15-01-1 p
1. --- మన ఆహారానికి కావాల్సిన సరుకులు మనకు అందించే వ్యాపారుల్లో చాలామంది
\VUgras{కప్పు ఉన్న సంత} \VUgras{మంటపం} \textsuperscript{(1)} కింద, లేదా పక్క వీధుల్లో
గుమిగూడతారు.

%%K t15-02-1 p
2. --- \VUgras{పంది మాంసం అమ్మేవాడు} \textsuperscript{(2)} — అతను \VUgras{హామ్}
\textsuperscript{(3)}, \VUgras{రక్తపు సాసేజీ} \textsuperscript{(4)}, \VUgras{సాసేజీ}
\textsuperscript{(5)}, \VUgras{పేగుల సాసేజీ} \textsuperscript{(6)},
\VUgras{పంది కొవ్వు మాంసం} \textsuperscript{(7)} అమ్ముతాడు — \VUgras{వెన్న, జున్ను అమ్మేది}
\textsuperscript{(8)} పక్కనే నిలబడతాడు; ఆమె బల్లపై ఎర్రని తొక్కతో ఉన్న
\VUgras{హాలండు జున్నులు} \textsuperscript{(9)}, \VUgras{కామంబేర్ జున్నులు}
\textsuperscript{(10)}, \VUgras{గ్రుయేర్ చక్రాలు} \textsuperscript{(11)}, తాజా
\VUgras{వెన్న ముద్దలు} \textsuperscript{(12)} నిండి ఉన్నాయి.

%%K t15-03-1 p
3. --- ఆమె ముందు \VUgras{కూరగాయలు అమ్మేది} \textsuperscript{(13)} తన ఖాతాదారులకు అందమైన
\VUgras{టమాటాలు} \textsuperscript{(14)}, \VUgras{కాలిఫ్లవర్లు} \textsuperscript{(15)},
\VUgras{సలాడు ఆకుకూర} \textsuperscript{(16)}, \VUgras{క్యాబేజీలు} \textsuperscript{(17)},
\VUgras{గుమ్మడి కాయలు} \textsuperscript{(19)}, \VUgras{కర్బూజాలు} \textsuperscript{(18)},
చివరికి \VUgras{పుట్టగొడుగులు} \textsuperscript{(20)} కూడా చూపిస్తుంది. కానీ
\VUgras{బీరు సరఫరా కుర్రాడు} తన \VUgras{సరుకు బండిని} \textsuperscript{(21)} —
\VUgras{బీరు పీపాలతో} \textsuperscript{(22)} నింపినది — సరిగ్గా ఆమె బల్ల ముందే ఆపడంతో ఆమె
ఖాతాదారులు దగ్గరికి రాలేకపోతున్నారు. తోటపని చేసే స్త్రీ ఒకామె ఆమెకు
\VUgras{ఆర్టిచోకుల} \textsuperscript{(23)} బుట్ట తెచ్చిపెడుతోంది.

%%K t15-tit-2 sub
\VUsaut{4.0mm}
\VUpk{10.2pt}{40}{అమ్మకం, కొనుగోలు.}
\VUsaut{3.0mm}

%%K t15-04-1 p
4. --- సంతలో సందడి ఎక్కువ, ఎందుకంటే \VUgras{వేలం అమ్మకం} \textsuperscript{(24)} వేళ
వచ్చేసింది. కొనుగోళ్ళకు వచ్చే \VUgras{గృహిణులు} \textsuperscript{(26)}
\VUgras{పేగులు అమ్మేది} \textsuperscript{(25)} దుకాణం ముందు ఆగుతారు —
\VUgras{పొట్ట పేగులు} గానీ \VUgras{దూడ తలకాయ} గానీ కొనేందుకు.

%%K t15-05-1 p
5. --- లావుపాటి \VUgras{వంటవాడు} \textsuperscript{(27)} ఒకడు చాలా మంచి \VUgras{ట్యూనా చేప}
కోసం బేరం చేస్తున్నాడు — \VUgras{చేపలు అమ్మేది} \textsuperscript{(28)} దగ్గర; ఆమె తన ఇష్టం
వచ్చినట్టు ధరలు పెంచుతుంది. ఆమెకు \VUgras{మలుగు చేపలు}, \VUgras{నాలుక చేపలు},
\VUgras{ట్రౌట్ చేపలు}, \VUgras{బొచ్చె చేపలు} పెద్ద మొత్తంలో వచ్చాయి. ఆమె
\VUgras{ఉప్పు చేపా} \VUgras{ఆల్చిప్పలూ} చాలా తాజాగా ఉన్నాయి.

%%K t15-tit-3 sub
\VUsaut{4.0mm}
\VUpk{10.2pt}{40}{చౌక సరుకులు, ఖరీదైన సరుకులు.}
\VUsaut{3.0mm}

%%K t15-06-1 p
6. --- ఆమె పక్కనే కూర్చున్న \VUgras{కోళ్ళు అమ్మేది} \textsuperscript{(29)} చాలా
నిజాయితీపరురాలు. \VUgras{టర్కీ కోడి}, \VUgras{గూస్ పక్షి}, \VUgras{బాతు} గానీ మామూలు
\VUgras{కోడిపిల్ల} గానీ అమ్మేటప్పుడు ఆమె సరైన ధరే అడుగుతుంది.

%%K t15-07-1 p
7. --- చౌక మూలలో, మొదటి అంతస్తులో, కొంతకాలంగా \VUgras{గుడ్డల వ్యాపారి}
\textsuperscript{(30)} ఒకడు కూర్చున్నాడు; అతని దగ్గర కొనేవాళ్ళకు అందమైన
\VUgras{బల్లగుడ్డలు}, \VUgras{చేతిగుడ్డలు}, \VUgras{ఏప్రాన్లు},
\VUgras{తుడుపు గుడ్డలు} ఎప్పుడూ దొరుకుతాయి. అదే అంతస్తులో \VUgras{కత్తుల పనివాడు}
\textsuperscript{(31)} తన పనిశాల పెట్టుకున్నాడు. సంతలోని అమ్మేవాళ్ళ \VUgras{కత్తులకు}
అతను పదును పెడతాడు, తోటపని చేసేవాళ్ళకు అన్నిటికంటే గట్టి \VUgras{ఉక్కుతో}
\VUgras{కొడవళ్ళు} చేస్తాడు.

%%K t15-tit-4 sub
\VUsaut{4.0mm}
\VUpk{10.2pt}{40}{కిరాణా సరుకులు.}
\VUsaut{3.0mm}

%%K t15-08-1 p
8. --- అతని పనిశాల కింద, కొంతకాలంగా, పెద్ద ఆహార సరుకుల దుకాణం తెరిచారు. అది ఇక పాతకాలపు
సాదా, చిన్న \VUgras{కిరాణా దుకాణం} \textsuperscript{(32)} కాదు — ఆ దుకాణం రోజువారీ సరుకులే
అమ్మేది : \VUgras{టీ} \textsuperscript{(33)}, \VUgras{చాక్లెట్టు} \textsuperscript{(34)},
\VUgras{చక్కెర గడ్డ} \textsuperscript{(37)}, \VUgras{సబ్బు} \textsuperscript{(38)}. ఇక్కడ
మాత్రం ఏ సరుకైనా అమ్ముతారు — \VUgras{కరకరలాడే బిస్కెట్లు},
\VUgras{డబ్బా ఆహారం} \textsuperscript{(35)}, \VUgras{మద్యాలు} \textsuperscript{(36)},
\VUgras{రమ్ము}, \VUgras{కోన్యాక్}, \VUgras{కిర్ష్}, \VUgras{సోంపు మద్యం},
\VUgras{కురాసావో}. ఇక్కడ వేట మాంసం కూడా దొరుకుతుంది : \VUgras{కుందేలు}
\textsuperscript{(39)}, \VUgras{త్రష్ పక్షులు} \textsuperscript{(40)}, \VUgras{కౌజు పిట్టలు}
\textsuperscript{(41)}, \VUgras{పూరేడు పిట్టలు} \textsuperscript{(42)},
\VUgras{అడవి బాతు} \textsuperscript{(43)}, \VUgras{ఫీజెంట్} \textsuperscript{(44)} — ఇంకా
\VUgras{అరటిపండ్లు} \textsuperscript{(46)}, \VUgras{అనాసపండ్ల} వంటి విదేశీ పండ్లూ అంతే
సులభంగా.

%%K t15-09-1 p
9. --- కిరాణా \VUgras{కుర్రాడు} \textsuperscript{(45)} చాలా సాయపడేవాడు; అడిగినది ఇచ్చేందుకు
సిద్ధంగా ఉంటాడు : \VUgras{జామ్} \textsuperscript{(47)} నాలుగో వంతు గానీ అయిదో వంతు గానీ,
\VUgras{కొవ్వు} \textsuperscript{(48)}, \VUgras{ఊరగాయ దోసకాయలు} \textsuperscript{(49)},
ఉప్పులో ఊరిన \VUgras{సార్డిన్ చేపలు} \textsuperscript{(50)}.

%%K t15-09-2 p
\VUgras{ఖాతాదారులకు} \textsuperscript{(51)} ఇది చాలా అనుకూలం — రోజువారీ సరుకులతో పాటే
అరుదైన తొలిపంటలూ, రుచికరమైన పండ్లూ వాళ్ళకు ఇక్కడే దొరుకుతాయి : రసంతో నిండిన
\VUgras{బేరిపండ్లు} \textsuperscript{(52)}, \VUgras{రేగుపండ్లు} \textsuperscript{(53)},
\VUgras{ద్రాక్షలు} \textsuperscript{(54)}, \VUgras{పీచుపండ్లు} \textsuperscript{(55)},
\VUgras{బాదంపప్పులు} \textsuperscript{(56)}, \VUgras{అక్రోట్లు} \textsuperscript{(57)}.

%%K t15-tit-5 sub
\VUsaut{4.0mm}
\VUpk{10.2pt}{40}{ఖాతాదారులు.}
\VUsaut{3.0mm}

%%K t15-10-1 p
10. --- \VUgras{బియ్యం} \textsuperscript{(58)}, \VUgras{మసూరు పప్పు} \textsuperscript{(59)}
పొగబెట్టిన \VUgras{హెర్రింగు చేపల} \textsuperscript{(60)} పెట్టె పక్కనే ఉన్నాయి;
\VUgras{బంగాళదుంపా} \textsuperscript{(61)} \VUgras{చిక్కుడు గింజలూ} \textsuperscript{(63)}
\VUgras{ఆపిల్స్} \textsuperscript{(62)}, \VUgras{అంజూర పండ్లు} \textsuperscript{(64)},
\VUgras{ఎండు రేగుపండ్లు} \textsuperscript{(65)}, \VUgras{హేజల్‌నట్ల} \textsuperscript{(66)}
పక్కన ఉన్నాయి. ఆ దుకాణంలో తాజా \VUgras{గుడ్లూ} \textsuperscript{(67)}
\VUgras{ఆలివ్‌లూ} \textsuperscript{(68)} దొరుకుతాయి; అందుకే \VUgras{బుట్టలూ}
\textsuperscript{(70)} \VUgras{సరుకుల వలసంచులూ} \textsuperscript{(73)} చాలా వేగంగా
నిండిపోతాయి.

%%K t15-11-1 p
11. --- ఆ చక్కటి దుకాణపు \VUgras{కాఫీ గింజలకు} \textsuperscript{(72)}
\VUgras{పనిమనుషుల} \textsuperscript{(69)}, \VUgras{వంట స్త్రీల} మధ్య తగిన పేరు వచ్చింది —
ఎందుకంటే ముసలి \VUgras{కిరాణా వ్యాపారే} \textsuperscript{(71)} వాటిని ఎంతో జాగ్రత్తగా
వేయిస్తాడు.

%%K t15-tit-6 sub
\VUsaut{4.0mm}
\VUpk{10.2pt}{40}{చూడదగ్గవి.}
\VUsaut{3.0mm}

%%K t15-12-1 p
12. --- అందరూ సరుకులు కొనేందుకే సంతకు రారు. మంటపానికి వెళ్ళే సందులో రమణీయమైన ఆ రద్దీలో
రకరకాల మనుషులు కనిపిస్తారు. మర్యాదస్తుడైన \VUgras{మాంసం దుకాణపు కుర్రాడి}
\textsuperscript{(75)} పక్కనే — అతను తన ముందు \VUgras{తోపుడు బండిని} \textsuperscript{(74)}
తోసుకెళ్తున్నాడు — ఇదిగో సెలవుపై వచ్చిన ఇద్దరు సైనికులు, తమ మెరిసే దుస్తులు
చూపించుకుంటూ గర్వంగా : నీలి \VUgras{దోల్మాన్ కోటుతో}, \VUgras{ఈకల టోపీతో}
\VUgras{హుస్సారు అశ్వసైనికుడు} \textsuperscript{(76)}; ఉబ్బిన చేతుల చొక్కాతో, తలపై
\VUgras{ఫెజ్ టోపీతో} \VUgras{జువావ్ దళపతి}.

%%K t15-13-1 p
13. --- \VUgras{రొట్టెలు చేసేవాడు} \textsuperscript{(80)} — అతను \VUgras{రొట్టెల దుకాణం}
\textsuperscript{(78)} గడపమీద నిలబడ్డాడు — ఇప్పుడే తన \VUgras{రొట్టె} \textsuperscript{(79)}
పిండి కలిపి, పొయ్యిలోంచి \VUgras{క్రొసాంలు} \textsuperscript{(81)},
\VUgras{చిన్న రొట్టెలు} \textsuperscript{(82)}, \VUgras{గుండ్రటి రొట్టెలు}
\textsuperscript{(83)} తీశాడు; అవి ఇప్పుడు అద్దాల అరలో ఉన్నాయి.

%%K t15-14-1 p
14. --- రొట్టెల దుకాణం పైన నేర్పరి \VUgras{నార బట్టలు కుట్టేది} \textsuperscript{(84)}
ఉంటుంది — ఆమె \VUgras{ఎంబ్రాయిడరీ చేసే స్త్రీలను} \textsuperscript{(85)} కొందరిని పనిలో
పెట్టుకుంది. ఆమె పక్కన \VUgras{బట్టలుతికి ఇస్త్రీ చేసేది} \textsuperscript{(86)}
ఉంటుంది; ఆమె \VUgras{నార బట్టలను} \textsuperscript{(88)} ఉతికి ఆరబెట్టాక గంజి పెట్టి,
\VUgras{ఇస్త్రీపెట్టెతో} \textsuperscript{(87)} నొక్కుతుంది.

%%K t15-tit-7 sub
\VUsaut{4.0mm}
\VUpk{10.2pt}{40}{మాంసం, చేపలు, కూరగాయలు.}
\VUsaut{3.0mm}

%%K t15-15-1 p
15. --- నేల అంతస్తులో ఉన్న \VUgras{మాంసపు దుకాణంలో} \textsuperscript{(89)} అన్ని రకాల
\VUgras{మాంసాలూ} అమ్ముతారు — \VUgras{గొర్రె మాంసం} \textsuperscript{(90)},
\VUgras{ఎద్దు మాంసం} \textsuperscript{(94)}, \VUgras{దూడ మాంసం} \textsuperscript{(92)} —,
\VUgras{గొర్రె తొడలు}, \VUgras{బీఫ్‌స్టేకులు}, \VUgras{కాల్చిన ముక్కలు},
\VUgras{కట్‌లెట్లు} రూపంలో. \VUgras{మాంసం కోసేవాడు} \textsuperscript{(93)} ఆ మాంసమంతా
ముక్కలు చేసి, \VUgras{మూలుగ ఎముకలను} \textsuperscript{(96)} విడిగా పెడతాడు. దుకాణం
తలుపు దగ్గర \VUgras{ఆయిస్టర్లు అమ్మేది} \textsuperscript{(97)} ఉంటుంది — ఆమె
\VUgras{ఆయిస్టర్లు} \textsuperscript{(98)} అమ్ముతుంది, కోరితే వాటిని తెరిచి ఇస్తుంది.

%%K t15-16-1 p
16. --- ఆ వ్యాపారులందరూ పన్నో చోటు రుసుమో కడతారు; అందుకే \VUgras{పోలీసు}
\textsuperscript{(100)} పాపం \VUgras{వీధిలో కూరగాయలు అమ్మే స్త్రీలను}
\textsuperscript{(99)} కనికరం లేకుండా తరుముతాడు — వాళ్ళు తమ చిన్న బళ్ళపై
\VUgras{కారెట్లు}, \VUgras{ముల్లంగి}, \VUgras{టర్నిప్పులు}, \VUgras{లీక్ కాడలు},
\VUgras{ఆస్పరాగస్ కట్టలు}, \VUgras{పచ్చి బఠానీలు}, \VUgras{పాలకూర},
\VUgras{చుక్కకూర}, \VUgras{వాటర్‌క్రెస్}, \VUgras{పార్స్లీ} తిప్పుతూ వాళ్ళకు పోటీ
ఇస్తారు.

%%K t15-tit-8 sub
\VUsaut{4.0mm}
\VUpk{10.2pt}{40}{పోలీసు జాగ్రత్త!}
\VUsaut{3.0mm}

%%K t15-17-1 p
17. --- \VUgras{వెల్లుల్లి} \textsuperscript{(101)}, \VUgras{ఉల్లిపాయలు} అమ్మే కుర్రాడికి
కూడా బాగా తెలుసు — తానూ సంత దగ్గర తన సరుకు అమ్మకూడదని. అతను పారిపోతాడు,
\VUgras{సముద్రపు ఎండ్రకాయలూ} \VUgras{లాబ్‌స్టర్లూ} \VUgras{అమ్మేవాడి}
\textsuperscript{(102)} \VUgras{పీతల}, \VUgras{ఏటి రొయ్యల}, \VUgras{రొయ్యల} బుట్టను
పడగొట్టే ప్రమాదంతో పరుగెత్తుతూ.

%%K t15-18-1 p
18. --- అందరూ కదులుతూ పెద్ద గోల చేస్తున్నారు; అయినా తిరుగుతూ అమ్మే
\VUgras{అద్దాలు వేసేవాడి} \textsuperscript{(103)} గొంతూ, \VUgras{బొగ్గు అమ్మేవాడి}
\textsuperscript{(104)} కేకా స్పష్టంగా వినిపిస్తాయి — ఆ బొగ్గువాడు
\VUgras{బొగ్గు మూట} \textsuperscript{(105)} అందించేందుకు తన బండిని కొన్ని క్షణాలు
వదిలివచ్చాడు.

\VUsaut{6.0mm}
\VUfilet{22mm}
